\documentclass[11pt]{article}
\usepackage[T1]{fontenc}
\usepackage{lmodern,amsmath,amssymb,amsthm,mathtools}
\usepackage[margin=1in]{geometry}
\usepackage{microtype,booktabs,array}
\usepackage[hidelinks]{hyperref}
\usepackage{url}
\newtheorem{theorem}{Theorem}
\newtheorem{lemma}[theorem]{Lemma}
\newtheorem{corollary}[theorem]{Recalled corollary}
\theoremstyle{remark}\newtheorem{remark}[theorem]{Remark}
\newcommand{\Q}{\mathbb Q}
\newcommand{\Z}{\mathbb Z}
\newcommand{\SL}{\mathrm{SL}}
\newcommand{\V}{\mathcal V}
\newcommand{\M}{\mathcal M}
\newcommand{\N}{\mathcal N}
\newcommand{\F}{\mathcal F}
\newcommand{\id}{\mathrm{id}}
\newcommand{\Tr}{\operatorname{Tr}}
\newcommand{\ad}{\operatorname{ad}}
\newcommand{\lk}{\operatorname{lk}}
\newcommand{\Diff}{\operatorname{Diff}}
\newcommand{\Int}{\operatorname{int}}
\newcommand{\sig}{\operatorname{symb}}
\setlength{\emergencystretch}{2em}
\title{Finite-type invariants on free-exterior string links:\\
low-degree restrictions and an order-seven example}
\author{Alec Kriebel\\Independent researcher\\
\small\href{https://orcid.org/0009-0001-9320-500X}{ORCID: 0009-0001-9320-500X}}
\date{13 October 2026 \quad Version 1.0}
\begin{document}
\maketitle
\begin{abstract}
We give applications of published finite-type classification and classical
genus-two symmetry to the locus of string links with free exterior group.
The category is ordinary unframed smooth ordered string links with fixed
boundary, upward orientation, and rational invariants. Freeness means that
the actual three-dimensional exterior group is abstractly $F_k$.
Recalling a filtered Habegger--Masbaum corollary, the intersection of the
entire Milnor algebra with the two-strand order-at-most-two space is
$\operatorname{span}_{\Q}\{1,\ell,\ell^2\}$.
Meilhan's classification and known free tangles give a component Conway
coefficient outside the Milnor-plus-vanishing span at $(k,n)=(2,2)$, and
give zero vanishing kernel for every $k\geq2$ and $0\leq n\leq2$.
An exact marked genus-two reversal argument, combined with the prior
Duzhin--Karev ordinary orientation detector, gives a nonzero kernel
invariant of exact order seven on two strands.
The seven-wheel is credited to Bar-Natan following Duzhin--Karev's journal
proof. Historical priority of these combined free-exterior applications
is undetermined; no firstness is claimed.
\end{abstract}
\noindent\textbf{Status and disclosure.}
This is an unrefereed application note prepared with extensive AI
assistance in solving, checking, sourcing, and writing. Its internal
mathematical and priority audits were performed by AI systems. They are
not conventional human peer review or formal proof certification.
The accompanying kit records the review history through 13 October 2026
and the current preparation state.
The author is responsible for the manuscript.
\smallskip

\noindent\textbf{License.} \href{https://creativecommons.org/licenses/by/4.0/}{CC BY 4.0}
for this authored manuscript and authored supporting records; cited
third-party source bodies and artwork are excluded.

\section{Category, question, and scope}
Let $C$ be the corner-rounded oriented cylinder $D^2\times[0,1]$, with
distinct marked points $p_1,\ldots,p_k$ in the interior of $D^2$.
An element of $\SL_k$ is an ordered smooth proper embedding of $k$ intervals,
with strand $i$ directed from $(p_i,0)$ to $(p_i,1)$, standard in product
endpoint collars. Equivalence is smooth ambient isotopy fixing the boundary;
isotopies can be taken fixed on a smaller boundary collar. Upward specifies
orientation, and imposes no height-monotonicity condition. No strand framing
is part of an object.

Let $\V_n(k)$ denote the $\Q$-valued ordinary Vassiliev invariants of order
at most $n$, including constants. Their extension to transverse double
points uses the positive-minus-negative crossing rule, at self-crossings
and mixed crossings alike. Write
\begin{align}
\F_k&=\{L\in\SL_k:\pi_1(E(L))\cong F_k\},\label{eq:free}\\
\N_n(k)&=\{v\in\V_n(k):v(L)=0\text{ for every }L\in\F_k\},\label{eq:kernel}\\
\M_n(k)&=\V_n(k)\cap\Q[\mu_I:I\text{ any multi-index}].\label{eq:milnor}
\end{align}
Here $E(L)$ is the compact actual three-dimensional exterior; the open
complement has the same fundamental group. $F_k$ is the abstract free
group of rank $k$. There is no assertion that the bottom meridians form
a free basis of that group. Neither nilpotent quotients nor a primitive
meridian condition replace \eqref{eq:free}. In \eqref{eq:milnor} every
based string-link Milnor invariant is allowed, with repeated indices and
arbitrary length, using preferred zero-self-linking longitudes. A polynomial
is finite, but its displayed generators need not have short length.

Stanford's Question 2.13 in Ohtsuki's problem collection
\cite[\S2.6, p.~408]{Ohtsuki} asks separately about the vanishing kernel
and whether the Milnor and vanishing spaces span. The results treated here
are the following restricted applications.
\begin{theorem}\label{thm:main}
In the category just specified:
\begin{enumerate}
\item With $\ell=\mu_{12}=\lk$, the prior filtered Milnor corollary is
\[
\M_2(2)=\operatorname{span}_{\Q}\{1,\ell,\ell^2\}.
\]
If $a_1$ is the Conway $z^2$ coefficient of the first individually closed
strand, then $a_1\notin\M_2(2)+\N_2(2)$. Thus the universal affirmative
spanning assertion fails at $(k,n)=(2,2)$.
\item $\N_n(k)=0$ for all $k\geq2$ and $0\leq n\leq2$.
\item $\N_7(2)$ contains a nonzero rational invariant of exact ordinary
order seven. Consequently $\N_n(2)\ne0$ for every $n\geq7$.
\end{enumerate}
\end{theorem}
The first formula is recalled from the filtered concordance theorem of
Habegger--Masbaum \cite[Corollary 6.4 and Theorem 15.1]{HM}, after the
ordinary specialization explained below. The low-degree applications use
Meilhan's complete classification \cite{Meilhan}, Nogueira's free-tangle
construction \cite{Nogueira14}, and the Kauffman--Lambropoulou rational
example \cite{KL,KLDNA}. The order-seven application uses classical
genus-two hyperelliptic machinery \cite{BleilerMoriah,RS,BM}, Hatcher's
relative-ball theorem \cite{HatcherSmale}, and the Duzhin--Karev detector
\cite{DKv5}. This is an assembly and exposition of those inputs, with the
marking and collar obligations made explicit.

This work does not determine $\N_n(2)$ for $3\leq n\leq6$, the least
positive kernel order, higher-strand higher-order kernels, or higher-order
spanning and general classification. Failure of spanning at order two
does not by itself imply failure at larger order bounds. No claim about
the current worldwide open status of those residual questions is made.

\section{Published order-two coordinates and the Milnor corollary}
For a knot $K$, let $c(K)$ be its normalized Conway $z^2$ coefficient.
For a string link put
\[
a_i(L)=c(\widehat{L_i}),\qquad
w_{ij}(L)=c(p(L_i\cup L_j))-a_i(L)-a_j(L),
\]
where an individual closure is formed after forgetting all other strands,
and $p$ joins the two bottom ends and the two top ends with fixed short
plat caps. These are Meilhan's component and plat invariants
\cite[\S2.1 and Proposition 2.1]{Meilhan}. In particular
$c(U)=0$ and $c(T)=1$ for either trefoil chirality.

Meilhan's Theorem 2.4, in this ordinary ordered category, gives every
$v\in\V_2(k)$ an expansion
\begin{equation}\label{eq:classification}
v=Q((\ell_{ij})_{i<j})+\sum_{i<j<h}d_{ijh}\mu_{ijh}
   +\sum_i b_i a_i+\sum_{i<j}e_{ij}w_{ij},\qquad \deg Q\leq2.
\end{equation}
Here $\mu_{ijh}$ is a based distinct-index length-three invariant in that
source's convention. $Q$ includes all constants and linear terms, all
squares, all linking products with a shared label, and all products with
four distinct labels. We use the actual spanning theorem for functions,
not merely an associated-graded statement. For $k=2$ it reads
\begin{equation}\label{eq:six}
v=A+B\ell+C\ell^2+D a_1+E a_2+F w_{12}.
\end{equation}

\begin{corollary}[Habegger--Masbaum, ordinary degree-two specialization]
\label{cor:milnor}
The entire-algebra intersection $\M_2(2)$ equals
$\operatorname{span}_{\Q}\{1,\ell,\ell^2\}$.
\end{corollary}
\begin{proof}
Every based Milnor invariant, including repeated indices, is invariant
under smooth fixed-endpoint concordance
\cite[\S2.1.4 and \S5.1]{MY09}, where the attribution is to Casson
\cite{Casson}. Therefore every finite polynomial in arbitrary-length
Milnor invariants is concordance invariant. This property concerns the
resulting function and remains valid if high-degree terms in a presentation
cancel.

Pull an ordinary concordance invariant back to framed string links by
forgetting the framing. Habegger--Masbaum's tree quotient, Corollary 6.4,
and filtered Theorem 15.1 \cite[\S4 and \S15]{HM} control its order-at-most-two
symbols. In degree two the only connected tree has three legs. With only
two colors, two legs share a color; exchanging these changes the vertex
orientation, so antisymmetry kills that diagram over $\Q$. Framing
independence kills the two self-struts. Through degree two the surviving
ordinary tree diagrams are therefore the empty diagram, the mixed strut,
and its square. They are integrated by $1,\ell,\ell^2$; the filtered theorem
gives the upper bound of three dimensions. Pure two-braids attain every
integer value of $\ell$, proving independence. These three functions belong
to the Milnor algebra and have orders at most $0,1,2$, respectively. The
upper and lower inclusions follow.
\end{proof}

\subsection{An elementary alternative: the full concordance equations}
This alternative proof supplies explicit ordinary tests, without claiming
a new Milnor theorem. Let $S=T\#(-T)$ be a long square knot, where $-T$ is
the reversed mirror. The standard smooth concordance inverse gives a
null-concordance fixed at the long endpoints
\cite[pp.~1--2]{CochranHarvey}. Conway multiplicativity and the mirror and
reversal identities give
\begin{equation}\label{eq:square}
\nabla_{K\#J}(z)=\nabla_K(z)\nabla_J(z),\quad
\nabla_{\bar K}(z)=\nabla_K(-z),\quad
\nabla_{K^r}(z)=\nabla_K(z),\quad c(S)=2.
\end{equation}
The first identity is recalled in \cite[\S1]{Conant}; the other two follow
from the normalized skein characterization.

Insert $S$ in a small ball on the first or second strand of the identity
link. The resulting $J_1,J_2$ are null-concordant and have linking zero;
their $(a_1,a_2,w)$ rows are $(2,0,0)$ and $(0,2,0)$, since each plat
contains that one local knot.

For a third test choose a product-collared null-concordance rectangle for
$S$ in $C\times[0,1]$. Its oriented rank-two normal bundle over the disk
is trivial. Prescribe the usual spatial separating normal direction on
the top edge and both endpoint edges. These edges form one contractible
boundary arc; the nonzero field extends over the rectangle and can be
chosen product-like in the boundary collars. The bottom direction is
determined by this extension and is not independently prescribed. Small
push-offs by $\pm\epsilon$ form disjoint fixed-endpoint concordance
rectangles with identity top. Their bottom link $C(S)$ has linking zero
by concordance, and each individual closure is $S$. The bottom field sweeps
out an embedded narrow ribbon rectangle whose short sides are the fixed
plat caps. Its boundary is the plat and bounds that embedded disk in $C$,
so the plat is unknot. Thus
\begin{equation}\label{eq:slice-rows}
\begin{array}{c|rrr}
 &a_1&a_2&w\\\hline
 J_1&2&0&0\\J_2&0&2&0\\C(S)&2&2&-4
\end{array}
\qquad \det=-16.
\end{equation}
For a concordance-invariant function in \eqref{eq:six}, subtracting its
identity value yields
\[
2D=0,\qquad 2E=0,\qquad 2D+2E-4F=0.
\]
Over $\Q$ these kill every non-linking coefficient, giving the same
corollary. No free-exterior assertion is made for the three auxiliary
links. This argument also explains why rational coefficients are essential;
it gives no arbitrary-coefficient-group theorem.

\section{Known free two-tangles and their ordinary markings}
\begin{lemma}\label{lem:witnesses}
There are $R,L,L^s\in\F_2$ with
\begin{equation}\label{eq:witnesses}
(a_1,a_2,w)(R)=(0,0,1),\qquad
(a_1,a_2)(L)=(1,5),\qquad(a_1,a_2)(L^s)=(5,1).
\end{equation}
No linking number or value of $w(L),w(L^s)$ is needed.
\end{lemma}
\begin{proof}
Nogueira's institutional preprint \cite[Appendix \S4, pp.~14--15,
Figures 9--10]{Nogueira14} gives a free two-arc ball tangle with individual
capped knots trefoil and $T(3,-4)$. Its explicit trefoil tunnel and tunnel-end
slide establish the actual handlebody exterior. This is the appendix
construction, not a consequence of the paper's main essential-surface
theorem. The related tunnel-slide method appears already in
\cite[\S8]{Nogueira13}. We use the full institutional drawings: the later
arXiv edition \cite{Nogueira17} has missing relevant artwork, and is not the
controlling visual source.

Round the finite tame PL corners in disjoint small neighborhoods and
straighten endpoint collars. This preserves complement topology and
individual knot types. Choose one endpoint of each arc as bottom, and use
an orientation-preserving sphere isotopy moving the four labeled endpoints
to their prescribed destinations. Extend through a ball collar and make
the strand germs standard. Orient each strand bottom to top. This produces
$L$ in the specified smooth ordered category, preserving its actual
handlebody exterior. The exterior boundary is the sphere with four disks
removed and two annuli attached, a closed genus-two surface; its group is
therefore $F_2$.

To compare individual caps, first discard the other strand. Any two simple
arcs joining the surviving two endpoints on the unpunctured sphere are
isotopic relative to those endpoints, by bigon removal. Extend across the
outside closing collar. Thus the individual knot types survive the chosen
marking. This is not a statement about arbitrary cap systems on a
four-punctured sphere. Willerton's normalized torus formula,
credited there to Alvarez--Labastida, is
\cite[\S1 and \S3, p.~5]{Willerton}
\[
c(T(p,q))=\frac{(p^2-1)(q^2-1)}{24}.
\]
It gives $c(T(3,-4))=5$ and $c(T)=1$. Mirror and reversal do not change this
even coefficient. For a separate copy exchange the two basepoints by an
orientation-preserving disk diffeomorphism and assign the carried strands
their exchanged labels. The resulting new ordered example is $L^s$.
It is not an identification of the two examples in the ordered category.

For $R$ use the prior rational tangle $[1]+1/[2]$ of fraction $3/2$ in
\cite[pp.~20--21, Figure 15]{KL} and
\cite[p.~30, Figure 20]{KLDNA}. It has vertical connectivity
NW--SW and NE--SE: $[1]$ pairs diagonally and $1/[2]$ pairs vertically;
tracing the horizontal sum gives the two stated arcs and no circle.
This agrees with \cite[Theorem 6]{KL}. Label the left and right vertical
arcs $1,2$. Carry SW, SE to bottom points and NW, NE to top points, preserving
both numerator cap arcs. Such a sphere marking exists by matching disjoint
disk neighborhoods of the cap arcs and their complementary annulus;
extend it over the ball and straighten collars.

Rationality is a ball-pair homeomorphism to the trivial two-tangle
\cite[\S2]{KL}, supplied here by the inverse finite endpoint twists.
Its exterior is pair-of-pants times an interval, a genus-two handlebody.
This rationality homeomorphism need not preserve the numerator caps.
Each individual arc is boundary-parallel, so its closure is unknot after
discarding the other strand. The distinct cap-preserving marking makes the
standard plat the cited trefoil numerator. Its coefficient is one.
This proves \eqref{eq:witnesses} with the actual exterior groups.
\end{proof}
This lemma is source-backed existential geometry. It does not give a
marked crossing list with independently recognized Wirtinger group or
numerical linking vector. The stated source inputs and conversion argument
are sufficient for the evaluations below.

\section{The spanning counterexample at two strands and order two}
\begin{proof}[Proof of Theorem \ref{thm:main}(1), spanning assertion]
Let $P_j$ be the pure two-braid with $j$ full twists, signed so
$\ell(P_j)=j$. A height-dependent disk isotopy trivializes its exterior as
a twice-punctured disk times an interval. Hence $P_j\in\F_2$. Its individual
strands are unknotted, so $a_1(P_j)=0$.

If $a_1=m+u$ with $m\in\M_2(2)$ and $u\in\N_2(2)$, the recalled corollary
gives $m=A+B\ell+C\ell^2$. Evaluation at $P_0,P_1,P_{-1}$ yields
\[
A=0,\qquad A+B+C=0,\qquad A-B+C=0.
\]
Therefore $m=0$. Evaluation at the free witness $L$ from
Lemma \ref{lem:witnesses} would then give $1=a_1(L)=u(L)=0$.
Equivalently, the evaluation matrix on columns $(1,\ell,\ell^2,a_1)$ is
\[
\begin{pmatrix}1&0&0&0\\1&1&1&0\\1&-1&1&0\\1&t&t^2&1\end{pmatrix},
\qquad\det=2\quad(t=\ell(L)).
\]
This proves the obstruction independently of the unknown marking-dependent
linking value. No matching of the full Milnor data to a braid is asserted.
\end{proof}

\section{Pure braid tests for the complete low-degree expansion}
Every pure $k$-braid has actual exterior diffeomorphic to a $k$-punctured
disk times an interval, so belongs to $\F_k$. Its individual strands are
unknotted. Its deletion to a pair is a pure two-braid whose two-plat is
unknot: consecutive twists slide across the endpoint cap, equivalently
this plat has one bridge. Thus
\begin{equation}\label{eq:braid-vanish}
a_i(P)=w_{ij}(P)=0\quad\hbox{for every pure braid }P.
\end{equation}

\subsection{An isolated triple and explicit leading terms}
On three strands take $A=\sigma_1^2$, $B=\sigma_2^2$ and
$\beta=ABA^{-1}B^{-1}$. Every two-strand deletion is identity, since it kills
one of $A,B$ and the other cancels its inverse. All linking numbers vanish.
Use the Artin convention
\[
\Phi_{\sigma_i}(x_i)=x_ix_{i+1}x_i^{-1},\qquad
\Phi_{\sigma_i}(x_{i+1})=x_i,\qquad
\Phi_{AB}=\Phi_A\circ\Phi_B.
\]
For a pure braid put $\Phi(x_j)=c_jx_jc_j^{-1}$ and normalize the $x_j$
exponent of $c_j$ to zero by right multiplication by an $x_j$ power.
These conjugators give the preferred longitudes, up to inversion of the
conjugation convention, by the parallel/Wirtinger interpretation
\cite[\S4.1 and \S5.2]{BNhomotopy}. Use the standard Magnus expansion
$x_i\mapsto1+X_i$, rather than the exponential variant in that source.
For $\beta$, direct substitution and free reduction give
\begin{align*}
c_1&=x_1x_2x_1^{-1}x_3x_1x_2^{-1}x_1^{-1}x_3^{-1},\\
c_2&=x_3^{-1}x_1^{-1}x_3x_1,\\
c_3&=x_3^{-1}x_1^{-1}x_3x_1x_2^{-1}x_1^{-1}
      x_3^{-1}x_1x_3x_1x_2x_1^{-1}.
\end{align*}
Using $E(x_i^{-1})=1-X_i+X_i^2+O(3)$, their expansions are
\begin{align}
E(c_1)&=1+X_2X_3-X_3X_2+O(3),\notag\\
E(c_2)&=1-X_1X_3+X_3X_1+O(3),\label{eq:magnus}\\
E(c_3)&=1+X_1X_2-X_2X_1+O(3).\notag
\end{align}
The selected triple coefficient is $\epsilon=\pm1$. All linear terms vanish,
so the possible lower-linking convention corrections vanish too.

For any chosen triple of labels, take a subdisk containing exactly those
points, perform this braid there, and keep the others vertical. Local
meridians are conjugates of the ambient meridians with the same linear
Magnus terms. A substitution by such conjugates preserves the degree-two
leading term in \eqref{eq:magnus}. For completeness an arbitrary change of
whisker introduces a conjugator of the form
\[
\Phi(q)c_jq^{-1}=(\Phi(q)q^{-1})(qc_jq^{-1}).
\]
Here $c_j\in\Gamma_2$ and $\Phi$ is identity modulo $\Gamma_3$, including on
inactive generators. Hence $\Phi(q)q^{-1}\in\Gamma_3$; the displayed
quadratic coefficient is unchanged. Deletion to any other triple leaves
at most two active strands and gives identity. Thus this pure braid isolates
one $\mu_{ijh}$ with value $\pm1$ and all other distinct-index triple values
zero, including when the chosen labels are nonadjacent.

\subsection{The full integer linking lattice}
A thin disk neighborhood of an arc joining any chosen pair of basepoints
and avoiding the others supports a full twist with linking one on that
pair and zero on all other pairs. Integer powers and products realize
every vector in $\Z^{\binom{k}{2}}$. This multiplication is used only within
pure braids.

If $v\in\N_2(k)$ is expanded as \eqref{eq:classification}, identity gives
$Q(0)=0$. The isolated triples and \eqref{eq:braid-vanish} give
$\epsilon d_{ijh}=0$. After all $d_{ijh}$ vanish, the full linking lattice
gives $Q(z)=0$ for every integral vector $z$. A rational polynomial
vanishing on that lattice is zero: fix all but one integral coordinate
and use the one-variable root theorem, then induct on the number of
variables. This proves the entire polynomial conclusion, rather than
extrapolating from a finite sample. We are left with
\begin{equation}\label{eq:remaining}
v=\sum_i b_i a_i+\sum_{i<j}e_{ij}w_{ij}.
\end{equation}

\section{Separated insertion and the all-strand zero kernel}
\begin{lemma}\label{lem:insertion}
Fix a pair $i,j$. A two-strand witness $S$ can be inserted in a subcylinder
separated from the other $k-2$ vertical strands by a disk. For this insertion
$J_{ij}(S)$,
\begin{equation}\label{eq:group-split}
\pi_1(E(J_{ij}(S)))\cong\pi_1(E(S))*F_{k-2}.
\end{equation}
Its active component and active plat values retain those of $S$; all
inactive component values and every other plat correction vanish.
\end{lemma}
\begin{proof}
Construct a proper base-disk arc $\delta$ cutting off exactly $p_i,p_j$:
use a thin regular neighborhood of a tree joining them to a boundary
interval, avoiding the finitely many other points. Both complementary
regions are disks. Insert $S$ in the selected region times $[0,1]$ with
its boundary collars, and keep the other strands vertical. The proper disk
$\Delta=\delta\times[0,1]$ misses the strands and their regular neighborhoods.
Cutting the compact exterior along $\Delta$ gives $E(S)$ and the trivial
$(k-2)$-string exterior. In collar-open neighborhoods their overlap retracts
to the disk $\Delta$. Van Kampen gives \eqref{eq:group-split}. For $k=2$ the
second factor is trivial.

Individual active closures are unchanged; inactive strands are unknotted.
For the active pair, delete the inactive strands before comparing caps.
In the unpunctured endpoint disks the transported caps isotope to the fixed
standard caps, so $w_{ij}$ retains its value. For a mixed active/outside pair,
choose each top and bottom cap to cross $\delta$ once, after other strands
have been discarded. Complete $\Delta$ through the outside closing ball
to a sphere meeting this plat exactly twice, once on each cap. Its two sides
give the two individual knot summands, perhaps reversing one orientation.
Thus the plat is their connected sum, and \eqref{eq:square} gives
\[
c(p(L_i\cup L_h))=a_i+a_h,\qquad w_{ih}=0.
\]
Outside/outside pairs are trivial. This proves every coordinate asserted.
\end{proof}
Only this disk-separated insertion preserves freeness here. No closure
property under arbitrary stacking or gluing is used, and deleting strands
is not asserted to preserve free exterior.

\begin{proof}[Proof of Theorem \ref{thm:main}(2)]
Take $v\in\N_2(k)$ and reduce it to \eqref{eq:remaining} by the braid tests.
For each pair $i<j$, the insertion $J_{ij}(R)$ is free by
\eqref{eq:group-split}, has all $a$ values zero, and has only $w_{ij}=1$
nonzero among the $w$ values. Thus $e_{ij}=0$. The free insertions
$J_{ij}(L)$ and $J_{ij}(L^s)$ then give
\[
b_i+5b_j=0,\qquad5b_i+b_j=0,\qquad
\det\begin{pmatrix}1&5\\5&1\end{pmatrix}=-24.
\]
Hence $b_i=b_j=0$ over $\Q$. Pairs $(1,j)$ cover all labels when $k\geq2$,
so $v=0$. Inclusion $\V_n(k)\subseteq\V_2(k)$ gives the conclusion for
$0\leq n\leq2$.
\end{proof}
For $k=2$, a finite full-rank certificate makes all unknown marking-dependent
values explicit. The rows on columns $(1,\ell,\ell^2,a_1,a_2,w)$ are
\begin{equation}\label{eq:six-matrix}
\begin{array}{c|rrrrrr}
P_0&1&0&0&0&0&0\\
P_1&1&1&1&0&0&0\\
P_{-1}&1&-1&1&0&0&0\\
R&1&r&r^2&0&0&1\\
L&1&s&s^2&1&5&u\\
L^s&1&t&t^2&5&1&v
\end{array}
\qquad\det=-48.
\end{equation}
The first three columns have upper block determinant $2$, and the bottom
three columns have block determinant $-24$, independently of $u,v$;
eliminating the lower linking entries gives their product. The identity
holds for all $r,s,t,u,v$, without assigning them numeric values.
The restriction $k\geq2$ matters: the cyclic-knot-group theorem makes the
one-strand free locus the unknot, on which a nonzero Conway coefficient
vanishes. That classical distinction is outside the all-pair argument.

\section{Exact marked reversal on every free two-string exterior}
Fix $p_1,p_2$ on the real diameter of the disk. A single fixed
orientation-preserving disk identification handles any other initial
basepoints. Define
\begin{equation}\label{eq:rho}
r(z,t)=(\bar z,1-t),\qquad(\rho\gamma)_i(s)=r(\gamma_i(1-s)).
\end{equation}
$r$ has three-dimensional determinant $+1$, and $\rho$ preserves labels,
upward orientations, and product endpoint collars. It is an involution.

\begin{lemma}\label{lem:reversal}
Every $L\in\F_2$ is ambient isotopic to $\rho L$ relative to the boundary,
with each label and orientation preserved. The isotopy can fix a sufficiently
small boundary collar.
\end{lemma}
\begin{proof}
We give the complete boundary-pattern and collar argument.

\emph{Recovering a handlebody.}
Delete the relative interiors of two standard tubes $N_i$ about the arcs.
The compact exterior $H$ has closed connected genus-two boundary, formed
from the outer four-holed sphere $P$ and the lateral annuli $A_i$.
It is irreducible: a sphere in $\Int H$ bounds its bounded-side ball in $C$
by smooth Schoenflies \cite[Theorem 1.1]{Hatcher3M}; every connected tube
meets $\partial C$ and misses the sphere, so lies outside that ball.
The ball is therefore contained in $H$.

A positive-genus closed surface group is not free, whereas every subgroup
of a free group is free. Thus a positive-genus boundary component cannot
inject into a free component group. The Loop Theorem
\cite[Theorem 3.1]{Hatcher3M} gives an essential compressing disk.
Cutting on it retains irreducibility: a sphere in a cut piece bounds a
ball in the uncut manifold, and the connected cutting disk, with boundary
on the exterior boundary and disjoint from that sphere, cannot enter this
ball. Reconstruction by a one-handle gives a free-product decomposition;
the cut component groups are free factors and remain free. Repeat.
The sum of $3g-2$ over positive-genus boundary components strictly decreases
under every essential compression (including torus and separating cases).
The final pieces have only spherical boundary and are irreducible, hence
are balls: a boundary sphere pushed into its collar bounds a ball on the
side away from that collar, and adjoining the collar gives the original
piece. Reverse the disk cuts to reconstruct a handlebody. Its boundary
genus is two. This uses actual group freeness, not a completion or special
meridian generating set.

\emph{Prescribing the entire closed-boundary map.}
Choose orientation-preserving tube charts
$q_i:D^2_\epsilon\times[0,1]\to N_i$, with
$q_i(0,s)=\gamma_i(s)$ and $q_i(z,s)=(p_i+z,s)$ near both ends.
The oriented normal bundle over an interval is trivial and the prescribed
end frames extend; integer twists are allowed construction choices.
Set
\begin{equation}\label{eq:tau}
\tau|_P=r|_P,\qquad
\tau(q_i(z,s))=q_i(\bar z,1-s)\quad(|z|=\epsilon).
\end{equation}
The formulas agree in full product neighborhoods of the seam circles.
Compatible equivariant corner rounding makes $\tau$ a smooth
orientation-preserving involution. It has six isolated fixed points:
two on $P$ and two on each $A_i$. The branched Euler formula is
$-2=2\chi(\partial H/\tau)-6$, giving a sphere quotient.

More precisely let $H_0$ be the trivial two-string exterior with the
same outer endpoint disks. The map $\phi:\partial H_0\to\partial H$ that
is literally identity on $P$ and uses $q_i$ on the annuli satisfies
\[
\tau=\phi(r|_{\partial H_0})\phi^{-1}.
\]
This is only a boundary identification; its extension has not been assumed.

\emph{Extension and exact boundary correction.}
The genus-two hyperelliptic class is central on the closed boundary and
extends over the standard handlebody
\cite[\S2, pp.~491--492]{RS}\cite[p.~272]{BM}.
Choose any orientation-preserving handlebody diffeomorphism $f:H\to H_0$.
Centrality gives
\[
[f\tau f^{-1}]=[r|_{\partial H_0}].
\]
Consequently $F_0=f^{-1}rf$ has boundary value $\sigma$ isotopic to $\tau$.
Take a smooth isotopy $\delta_t$ with $\delta_0=\id$ and
$\delta_1=\tau\sigma^{-1}$. In an inward collar $c(x,u)$, define
\[
K_0(c(x,u))=c(\delta_{\lambda(u)}(x),u),
\]
where $\lambda=1$ near the boundary and zero near the inner collar edge;
extend by identity. The unchanged $u$ coordinate shows this is a
diffeomorphism. Then $F=K_0F_0$ has $F|_{\partial H}=\tau$ pointwise.
Intermediate $\delta_t$ may move $P$ and the annuli, since no tubes are
glued at that stage. The endpoint is the prescribed map on the whole
closed boundary.

\emph{Correcting germs before smooth gluing.}
On a whole tube prescribe
$J_i=q_i\circ((z,s)\mapsto(\bar z,1-s))\circ q_i^{-1}$.
Together with $r$ near $P$ these define a compatible local involution $j$
near the boundary pattern, since the charts are standard at the seams.
Take an invariant local metric by averaging with $j^*$; its inward normal
collar obeys
\[
j(c(x,u))=c(\tau x,u).
\]
The extension $F$, already exact on the boundary, gives another collar
\[
d(x,u)=F(c(\tau^{-1}x,u)),\qquad d(x,0)=x.
\]
Uniqueness of smooth collars relative to their boundary, and isotopy
extension, yield a diffeomorphism $K_1$, supported near the boundary and
fixing it pointwise, such that $K_1(d(x,u))=c(x,u)$ for small $u$.
The inward transverse derivatives are positive, so the collars belong
to the same component. Therefore
\begin{equation}\label{eq:germ}
(K_1F)(c(x,u))=K_1(d(\tau x,u))=c(\tau x,u)=j(c(x,u)).
\end{equation}
Equality of boundary values alone would not have guaranteed smooth gluing;
\eqref{eq:germ} gives equality in a full neighborhood. Neither $K_1$
commuting with $\tau$ nor the global extension squaring to identity is needed.

Glue $K_1F$ to the $J_i$. The result is an orientation-preserving smooth
ball diffeomorphism $h$ satisfying, literally,
\[
h|_{\partial C}=r|_{\partial C},\qquad
h(\gamma_i(s))=\gamma_i(1-s),
\]
and equal to $r$ in a small boundary collar. Thus $g=rh$ is identity on
that collar and $g(\gamma_i(s))=(\rho\gamma)_i(s)$ for each label.
There is no residual boundary mapping class or pure-braid factor.

Hatcher's theorem gives contractibility of
$\Diff(D^3\,\mathrm{rel}\,\partial D^3)$
\cite[Appendix, p.~604, statement (1)]{HatcherSmale}, hence an isotopy
of $g$ to identity relative to the boundary. Parametric collar uniqueness
straightens this path to fix a uniformly smaller collar, with endpoint
corrections identity; compactness of the path parameter provides a common
positive width. This isotopy proves the lemma and extends by identity in
the long-link formulation. It is a smooth relative-ball theorem, not a
radial smooth Alexander-trick assertion.
\end{proof}
The classical genus-two machinery has older inversion and branched-cover
context in Bleiler--Moriah \cite[\S1.3]{BleilerMoriah}. The lemma is its
marked boundary/collar specialization with the explicit $r$ of
\eqref{eq:rho}. Centrality is specific to genus two; the argument supplies
no higher-strand reversal theorem. No separately prescribed framing is
retained.

\section{The prior signed seven-wheel and ordinary integration}
For the algebraic input we use the exact arXiv v5 body of Duzhin--Karev,
dated 25 July 2005 \cite{DKv5}. Its Lemma 4 gives leg parity under
simultaneous reversal; Proposition 2 gives the nonzero seven-wheel, and
Section 6, Lemma 8, gives its survival under ordinary deframing. The v1
framed argument alone would not suffice. Their later English journal proof
\cite[\S5, Remark 1]{DKjournal} explicitly credits Bar-Natan's earlier
discovery of the wheel in \texttt{table.m}; we carry that attribution.
The proof text inspected for this credit is author-uploaded and contains
production placeholders, so it is not certified as final publisher PDF
bytes. The journal bibliographic record and DOI are given in the references.

Let $H_7$ be the connected seven-wheel with cyclic leg colors $1121222$
and the vertex orientation of \cite[Proposition 2]{DKv5}. It has seven
trivalent and seven univalent vertices, so degree $(7+7)/2=7$. The colored
Jacobi reversal changes it by $(-1)^7=-1$. Under symmetrization $\chi$,
write $d=\chi(H_7)$ for the ordinary rational diagram class. The cited
deframing result says $d\ne0$ after imposing the ordinary 1T relations
that delete the two self-struts. Its qualification is important: a connected
nonsingle-strut class is primitive, and the two framing generators are
central degree-one primitives. The symmetrized wheel belongs to the
complementary ordinary subalgebra identified by Lemma 8. Connectedness
alone is not a general license to pass arbitrary framed classes through 1T.

The operation in \eqref{eq:rho} agrees with this reversal. For the standard
long lines $(i,0,t)$, reversing strand parameters is restored to upward
labeled lines by $(x,y,z)\mapsto(x,-y,-z)$; translation and rescaling give
exactly $r$. Ambient orientation is preserved, and reversing both tangents
at a crossing preserves its crossing sign. Thus $\rho$ preserves ordinary
finite-type order, with symbol involution $\tau_A$ satisfying
$\tau_A(d)=-d$. It is neither a label-exchanging half-turn nor a mirror.

\subsection{An optional exact finite trace certificate}
The following arithmetic supplements the published nonzero proof, and
does not replace its AS/IHX, deframing, or integration interfaces.
For matrices $X,Y\in\operatorname{Mat}_N(\Q)$ write $T_w$ for the trace of
the color word $w$, with color 1 assigned $X$, color 2 assigned $Y$,
and $T_\emptyset=N$. Since $\ad X=L_X-R_X$ and
$\Tr_{\operatorname{End}(\operatorname{Mat}_N)}(L_A R_B)=\Tr(A)\Tr(B)$,
the wheel trace has a signed expansion over all $2^7=128$ subsets of the
seven positions. Left factors retain order; right factors reverse order.
Cyclic grouping yields
\begin{equation}\label{eq:wheel}
P=N(T_{1121222}-T_{1122212})
   +3T_2(T_{112212}-T_{112122}).
\end{equation}
For the integer matrices
\[
X=\begin{pmatrix}0&1&1&1\\1&0&0&0\\0&-1&0&0\\-1&-1&-1&1\end{pmatrix},
\qquad
Y=\begin{pmatrix}0&0&-1&1\\-1&-1&1&1\\-1&0&0&0\\0&0&-1&-1\end{pmatrix},
\]
exact multiplication gives
\[
T_2=-2,\quad T_{1121222}=38,\quad T_{1122212}=20,
\quad T_{112212}=-12,\quad T_{112122}=-23.
\]
Thus $P=4(38-20)+3(-2)(-12+23)=6$. Reversing the word gives $-6$;
the antisymmetric difference is $12$. The Lie weight respects AS and IHX
by invariance of the trace pairing and the Jacobi identity, so this is
a rational nonzero diagram certificate, up to the fixed global vertex sign.
The new review checker separately derives all 128 terms and multiplies
independently constructed $16\times16$ adjoint matrices. Its frozen results
are supplementary stage evidence. The value $6$ is a diagram weight value,
not a value on a nonsingular string link.

\begin{proof}[Proof of Theorem \ref{thm:main}(3)]
The ordinary degree-seven rational diagram space is finite dimensional.
Choose a rational functional $W$ with $W(d)=1$, and put
\[
W^-=(W-W\circ\tau_A)/2.
\]
Since $\tau_A(d)=-d$, we have $W^-(d)=1$ and
$W^-\circ\tau_A=-W^-$. The rational ordinary integration theorem, in the
form used in \cite[\S2 and \S6]{DKv5}, supplies $f\in\V_7(2)$ with top
symbol $W^-$. Set
\begin{equation}\label{eq:v}
v=(f-f\circ\rho)/2.
\end{equation}
Sign-preserving naturality of singular resolution gives
\[
\sig_7(v)=\frac{W^- -W^-\circ\tau_A}{2}=W^-.
\]
This symbol is nonzero on $d$, so $v$ has exact order seven and is a
nonzero function. Indeed, if all nonsingular values were zero, every
singular alternating sum and the symbol would be zero.
For every $L\in\F_2$, Lemma \ref{lem:reversal} gives
$f(L)=f(\rho L)$, so $v(L)=0$. Hence $v\in\N_7(2)$, and the same member
belongs to $\N_n(2)$ for every $n\geq7$.
\end{proof}
The construction is existential: the choice of dual functional and
integration is not canonical. No named nonsingular link, numeric value on
such a link, or executable evaluation formula for $v$ is supplied.
The universal zero restriction follows from the geometric lemma; a finite
set of zero evaluations could not establish it. A minimum orientation
detecting degree, discussed in \cite[Remark 1.4]{MY14}, would not establish
a minimum kernel order, since a reversal-even invariant could lie in the
kernel. No minimum kernel claim is made here.

\section{Priority, evidence, and limitations}
The filtered Milnor corollary, complete order-two basis and plat invariant,
Nogueira free witnesses, rational $3/2$ example, genus-two machinery, and
Duzhin--Karev ordinary detector are published inputs, credited at their
points of use. The accompanying fresh priority report records a bounded
comparison of complete controlling bodies and relevant sections. It does
not establish originality of the combined free-exterior applications.

Two close older priority leads remain unread after recorded access failures:
Morimoto's 1995 Lemma 5 \cite{Morimoto} and Yannan Li's 2009 free-tangle paper
\cite{Li}. Their exact contents are unknown in this investigation. They
are priority unknowns, not mathematical dependencies of these proofs.
The complete Bleiler--Moriah article was recovered and read; it provides
older inversion context, with a different knot-exterior hypothesis.
No absence claim from an unread abstract or search excerpt supports this
note's conclusions. Further external scholarly input could improve the
priority investigation; only the human author may communicate externally.

The source/audit kit separates nineteen unchanged theoretical originals,
fresh mathematical derivations, new finite checks, bounded priority
comparison, and current preparation state. Historical acceptance records
retain their historical stages and are not package-readiness certificates.
The original edition delivered no executable checker; omitted historical
scripts and private fixtures were neither replayed nor simulated. New
review programs check finite word/Magnus, determinant, Alexander, and trace
arithmetic; these do not prove universal topology, published classification,
1T survival, or rational integration. Those are explicit mathematical
arguments and established external interfaces in this manuscript.
Copied third-party PDFs, PostScript, extracted text, OCR, and images are
excluded from the distributable authored kit. Bibliographic identities
and hashes are retained for source verification.

\begin{thebibliography}{99}\small
\bibitem{Ohtsuki} T. Ohtsuki (ed.), \emph{Problems on invariants of knots
and 3-manifolds}, Geometry \& Topology Monographs \textbf{4} (2002),
377--572; T. Stanford, Question 2.13, p.~408.
\url{https://msp.org/gtm/2002/04/gtm-2002-04-024s.pdf}.
\bibitem{HM} N. Habegger and G. Masbaum, \emph{The Kontsevich integral and
Milnor's invariants}, Topology \textbf{39} (2000), 1253--1289. Controlling
author preprint: 7 April 1999, \S4, Corollary 6.4, Theorem 15.1 and its
proof pp.~29--30. \url{https://webusers.imj-prg.fr/~gregor.masbaum/K39.ps.gz}.
\bibitem{Meilhan} J.-B. Meilhan, \emph{On Vassiliev invariants of order two
for string links}, J. Knot Theory Ramifications \textbf{14} (2005),
665--687. Exact preprint: math/0402036v2, 10 November 2004.
\url{https://arxiv.org/pdf/math/0402036v2}.
\bibitem{MY09} J.-B. Meilhan and A. Yasuhara, \emph{Characterization of finite
type string link invariants of degree $<5$}, Math. Proc. Cambridge Philos.
Soc. \textbf{148} (2010), 439--472. Exact preprint: 0904.1527v1,
9 April 2009. \url{https://arxiv.org/pdf/0904.1527v1}.
\bibitem{Casson} A. J. Casson, \emph{Link cobordism and Milnor's invariant},
Bull. London Math. Soc. \textbf{7} (1975), 39--40.
\url{https://doi.org/10.1112/blms/7.1.39}. Concordance attribution used
through \cite{HM,MY09}; original body not newly recovered.
\bibitem{CochranHarvey} T. D. Cochran and S. Harvey, \emph{The geometry of
the knot concordance space}, controlling author copy of arXiv:1404.5076v2,
11 July 2014, introduction pp.~1--2.
\url{https://math.rice.edu/~shelly/publications/arxiv_geometry_knot_conc.pdf}.
\bibitem{Conant} J. Conant, \emph{Chirality and the Conway polynomial},
Topology Proceedings \textbf{30} (2006), 153--162.
Controlling preprint: math/0503648v2, 22 February 2012, \S1.
\url{https://arxiv.org/pdf/math/0503648v2}.
\bibitem{Nogueira14} J. M. Nogueira, \emph{Knot complements with meridional
essential surfaces of arbitrarily high genus}, Coimbra preprint 14--09,
received 27 February 2014, Appendix \S4, pp.~14--15, Figures 9--10.
\url{https://www.mat.uc.pt/preprints/ps/p1409.pdf}.
\bibitem{Nogueira17} J. M. Nogueira, \emph{Prime knot complements with
meridional essential surfaces of arbitrarily high genus}, Topology Appl.
\textbf{194} (2015), 427--439, DOI 10.1016/j.topol.2015.05.088.
Later preprint: 1706.03719v1, with missing relevant appendix artwork.
\url{https://arxiv.org/pdf/1706.03719v1}.
\bibitem{Nogueira13} J. M. Nogueira, \emph{Tunnel number degeneration under
the connected sum of prime knots}, arXiv:1303.7275v2, 17 June 2013,
Definitions 1--3 and \S8.
\url{https://arxiv.org/pdf/1303.7275v2}.
\bibitem{KL} L. H. Kauffman and S. Lambropoulou, \emph{On the Classification
of Rational Knots}, exact preprint math/0212011v2, 27 November 2003,
pp.~6--9, 20--21, 39--41.
\url{https://arxiv.org/pdf/math/0212011v2}.
\bibitem{KLDNA} L. H. Kauffman and S. Lambropoulou, \emph{From Tangle
Fractions to DNA}, undated author exposition, p.~30, Figure 20.
\url{https://homepages.math.uic.edu/~kauffman/Dresden.pdf}.
\bibitem{Willerton} S. Willerton, \emph{On the first two Vassiliev
invariants}, arXiv:math/0104061v1, 5 April 2001, \S1 and \S3.
Torus formula credited there to Alvarez--Labastida.
\url{https://arxiv.org/pdf/math/0104061v1}.
\bibitem{BNhomotopy} D. Bar-Natan, \emph{Vassiliev homotopy string link
invariants}, J. Knot Theory Ramifications \textbf{4} (1995), 13--32;
author PDF edition 17 February 2015, \S4.1 and \S5.2.
\url{https://www.math.toronto.edu/drorbn/papers/homotopy/homotopy.pdf}.
\bibitem{DKv5} S. V. Duzhin and M. V. Karev, \emph{Detecting the orientation
of long links by finite type invariants}, exact arXiv:math/0507015v5,
25 July 2005, \S2, Lemma 4, Proposition 2, \S6 and Lemma 8.
\url{https://arxiv.org/pdf/math/0507015v5}.
\bibitem{DKjournal} S. V. Duzhin and M. V. Karev, \emph{Detecting the
orientation of string links by finite type invariants}, Funct. Anal. Appl.
\textbf{41}(3) (2007), 208--216.
\url{https://doi.org/10.1007/s10688-007-0018-7}.
Journal record: \url{https://www.mathnet.ru/eng/faa2865}.
The Bar-Natan discovery credit is in \S5, Remark 1 of the inspected
author-uploaded English journal proof:
\url{https://www.researchgate.net/publication/225483857}.
Its production placeholders prevent certification as final publisher bytes.
\bibitem{BleilerMoriah} S. A. Bleiler and Y. Moriah, \emph{Heegaard
splittings and branched coverings of $B^3$}, Math. Ann. \textbf{281}
(1988), 531--543. \url{https://doi.org/10.1007/BF01456837}.
Original article scans indexed at \url{https://eudml.org/doc/164437}.
\bibitem{RS} J. H. Rubinstein and M. Scharlemann, \emph{Genus two Heegaard
splittings of orientable three-manifolds}, Geometry \& Topology Monographs
\textbf{2} (1999), 489--553, \S2.
\url{https://msp.org/gtm/1999/02/gtm-1999-02-024p.pdf}.
\bibitem{BM} A. Bruno and M. Mecchia, \emph{On quotient orbifolds of
hyperbolic 3-manifolds of genus two}, Rend. Istit. Mat. Univ. Trieste
\textbf{46} (2014), 271--299, p.~272.
\url{https://rendiconti.dmi.units.it/volumi/46/014.pdf}.
\bibitem{Hatcher3M} A. Hatcher, \emph{Notes on Basic 3-Manifold Topology},
Theorems 1.1 and 3.1, pp.~1 and 56.
\url{https://pi.math.cornell.edu/~hatcher/3M/3Mfds.pdf}.
\bibitem{HatcherSmale} A. E. Hatcher, \emph{A proof of the Smale Conjecture,
$\Diff(S^3)\simeq O(4)$}, Ann. of Math. \textbf{117} (1983), 553--607,
Appendix p.~604, statement (1).
\url{https://pi.math.cornell.edu/~hatcher/Papers/SmaleConjecture.pdf}.
\bibitem{MY14} J.-B. Meilhan and A. Yasuhara, \emph{Abelian quotients of
the string link monoid}, Algebr. Geom. Topol. \textbf{14} (2014),
1461--1488, Remark 1.4 and \S3.2.
\url{https://msp.org/agt/2014/14-3/agt-v14-n3-p09-s.pdf}.
\bibitem{Morimoto} K. Morimoto, \emph{There are knots whose tunnel numbers
go down under connected sum}, Proc. Amer. Math. Soc. \textbf{123} (1995),
3527--3532. \url{https://doi.org/10.1090/S0002-9939-1995-1317043-4}.
Priority lead: Lemma 5 and full body unread; not a mathematical dependency.
\bibitem{Li} Y. Li, \emph{2-string free tangles and incompressible surfaces},
J. Knot Theory Ramifications \textbf{18} (2009), 1081--1087.
\url{https://doi.org/10.1142/S021821650900735X}.
Priority lead: full body unread; not a mathematical dependency.
\end{thebibliography}
\end{document}

% administrative fault control
